%%%%%%%%%%%%%%%%%%%%%%%%%%%%%%%%%%%%%%%%%%%%%%%%%%%%%%%%%%%%%%%%%%%%%%
% Lớp 10 | Véctơ (chưa xét tọa độ) | Tích của một số với véctơ
% Dạng 0H5?3-5: Phân tích 1 véctơ theo 2 véctơ không cùng phương
% Số câu: 9  (sinh tự động bởi tach_cau.py; từ nay thầy sửa trực tiếp file này)
%%%%%%%%%%%%%%%%%%%%%%%%%%%%%%%%%%%%%%%%%%%%%%%%%%%%%%%%%%%%%%%%%%%%%%

%%%%%%%%%%%%%%%% TRẮC NGHIỆM 4 PHƯƠNG ÁN %%%%%%%%%%%%%%%%

% nguồn: MaHoaNH_TracNghiem_Lop10
\begin{ex}%[0H5H3-5]%[ma:6379812f]
	Cho tam giác $ABC$ gọi $D$ là điểm sao cho $\overrightarrow{BD}=-\dfrac{2}{3}\overrightarrow{CB}$ và $M$ là trung điểm  $AD$, $\overrightarrow{BM}$ được phân tích theo hai vectơ $\overrightarrow{BA}$ và $\overrightarrow{BC}$. Hãy chọn khẳng định đúng trong các khẳng định sau:
	\choice 
		{$\dfrac{1}{2}\overrightarrow{BA}+\dfrac{2}{3}\overrightarrow{BC}$} 
		 {$\overrightarrow{BA}+\dfrac{2}{3}\overrightarrow{BC}$}
		{$\dfrac{1}{2}\left(\overrightarrow{BA}-\dfrac{2}{3}\overrightarrow{BC}\right)$} 
		{\True  $\dfrac{1}{2}\left(\overrightarrow{BA}+\dfrac{2}{3}\overrightarrow{BC}\right)$}
	\loigiai{ 
	Theo  quy tắc trung điểm ta có $\overrightarrow{BM}=\dfrac{1}{2}\left(\overrightarrow{BA}+\overrightarrow{BD}\right)=\dfrac{1}{2}\left(\overrightarrow{BA}-\dfrac{2}{3}\overrightarrow{CB}\right)=\dfrac{1}{2}\left(\overrightarrow{BA}+\dfrac{2}{3}\overrightarrow{BC}\right)$.}
\end{ex}

% nguồn: MaHoaNH_TracNghiem_Lop10
\begin{ex}%[0H5H3-5]%[ma:ffe023c1]
	\immini[thm]{ 
	Cho tam giác $ABC$ có trung tuyến $AM$, khẳng  định nào sau đây đúng?
	\choice 
		 {$\overrightarrow{AM}=\overrightarrow{AB}+2\overrightarrow{BM}$}
		{\True  $\overrightarrow{AM}=\dfrac{1}{2} \left(\overrightarrow{AB}+\overrightarrow{AC}\right)$}
		{$\overrightarrow{AM}=-\dfrac{1}{2} \left(\overrightarrow{AB}+\overrightarrow{AC}\right)$} 
		 {$\overrightarrow{AM}=\dfrac{1}{2} \left(\overrightarrow{AB}-\overrightarrow{AC}\right)$}
		 }{
		\begin{tikzpicture}[line join=round,line cap=round,thick,scale= 0.7]
				\coordinate (A) at (0,3); 
				 \coordinate (B) at (-2,0);
				\coordinate  (C) at (3,0);
				\coordinate  (M) at ($(B)!1/2!(C)$);
				\draw (A)--(B)--(C)--cycle; 
				 \draw (A)--(M);
				\foreach  \i/\g in {A/90,B/-90,C/-90,M/-90}{\draw[fill=white](\i) circle (1pt)
					($(\i)+(\g:3mm)$) node[scale=0.7]{$\i$};} 
			\end{tikzpicture}
		 }
	 \loigiai{
		Theo tính chất trung  điểm đoạn thẳng ta có $\overrightarrow{AB}+\overrightarrow{AC}=2\overrightarrow{AM}$.\\
		 Suy ra $\overrightarrow{AM}=\dfrac{1}{2} \left(\overrightarrow{AB}+\overrightarrow{AC}\right)$.
	} 
\end{ex}


%%%%%%%%%%%%%%%% ĐÚNG/SAI %%%%%%%%%%%%%%%%

% nguồn: MaHoaNH_DungSai_Lop10
\begin{ex}%[0H5H3-5]%[ma:6f7f88a6]
	Cho tam giác $ABC$ có $G$ là trọng tâm và $I$ là trung điểm của đoạn thẳng $BC$. Đặt $\overrightarrow{AB}=\vec{a}$, $\overrightarrow{AC}=\vec{b}$. Xét tính đúng  sai của các mệnh đề sau
	 \choiceTF
		 {\True $\overrightarrow{AI}=\dfrac{1}{2}\vec{a}+\dfrac{1}{2}\vec{b}$}
		 {$\overrightarrow{IG}=-\dfrac{1}{6}\vec{a}+\dfrac{1}{6}\vec{b}$}
		 {\True $\overrightarrow{BI}=-\dfrac{1}{2}\vec{a}+\dfrac{1}{2}\vec{b}$}
		{\True  $\overrightarrow{CI}=\dfrac{1}{2}\vec{a}-\dfrac{1}{2}\vec{b}$}
	\loigiai{ 
		\begin{itemchoice}
			\itemch Vì $I$  là trung điểm của $BC$ nên $\overrightarrow{AI}=\dfrac{1}{2}\vec{a}+\dfrac{1}{2}\vec{b}$.
			 \itemch Vì $\overrightarrow{IG}=-\dfrac{1}{3}\overrightarrow{AI}=-\dfrac{1}{6}\vec{a}-\dfrac{1}{6}\vec{b}$.
			 \itemch $\overrightarrow{BI}=\overrightarrow{AI}-\overrightarrow{AB}=-\dfrac{1}{2}\vec{a}+\dfrac{1}{2}\vec{b}$.
			\itemch $\overrightarrow{CI}=-\overrightarrow{BI}$. 
		\end{itemchoice}
	} 
\end{ex}


%%%%%%%%%%%%%%%% TRẢ LỜI NGẮN %%%%%%%%%%%%%%%%

% nguồn: MaHoaNH_TraLoiNgan_Lop10
\begin{bt}%[0H5H3-5]%[ma:80cb105e]
	Cho ba điểm $A$, $B$, $C$ không thẳng hàng. Gọi điểm $M$ là trọng tâm $\triangle ABC$ và thỏa mãn đẳng thức $\overrightarrow{MA}=\alpha \cdot\overrightarrow{MB}+\beta\cdot \overrightarrow{MC}~(\alpha, \beta \in \mathbb{R})$.  Tính $2\alpha+\beta$.
	 \shortans{3}
	\loigiai{ 
		Vì  $M$ là trọng tâm $\triangle ABC$ nên $\overrightarrow{MA}+\overrightarrow{MB}+\overrightarrow{MC}=\overrightarrow{0}\Leftrightarrow \overrightarrow{MA}=-\overrightarrow{MB}-\overrightarrow{MC}$.\\
		So sánh  với $\overrightarrow{MA}=\alpha\cdot \overrightarrow{MB}+\beta \cdot\overrightarrow{MC} \Rightarrow \alpha=-1$; $\beta=-1$ suy ra $2\alpha+\beta=-3$.
	} 
\end{bt}

% nguồn: MaHoaNH_TraLoiNgan_Lop10
\begin{bt}%[0H5H3-5]%[ma:7ffc77c0]
	Cho tam giác $ABC$.  Gọi $M$ là một điểm trên cạnh $BC$ sao cho $MB=2MC$. Nếu \break$\overrightarrow{AM} = m\overrightarrow{AB} + n \overrightarrow{AC}$ thì giá trị của $3m - 3n$ bằng bao nhiêu?
	 \shortans{-1}
	 \loigiai{
	\immini{ 
		Ta  có
		\[ 
		\overrightarrow{AM}=\overrightarrow{AC}+\overrightarrow{C M}=\overrightarrow{AC}-\dfrac{1}{3} \overrightarrow{BC} =\overrightarrow{AC}-\dfrac{1}{3}\left(\overrightarrow{AC}-\overrightarrow{AB}\right)=\dfrac{1}{3}  \overrightarrow{AB}+\dfrac{2}{3} \overrightarrow{AC}.
		\] 
		Suy ra $m = \dfrac{1}{3}$,  $n = \dfrac{2}{3}\Rightarrow 3m -3n = -1$.
		 }{
		\begin{tikzpicture}[scale=1,font=\footnotesize,line join=round,line cap=round,>=stealth]
				\path (0,0) coordinate  (B)++(0:3) coordinate (C)
					($(B)+(60:2)$)  coordinate (A)
					 ($(B)!2/3!(C)$) coordinate (M);
				\draw (A)--(B)--(C)--(A) 
					(A)--(M); 
				\foreach \x/\g in {A/90,B/-150,C/0,M/-90} \fill (\x) circle  (1pt)+(\g:.3) node {$\x$};
			\end{tikzpicture}
		} 
	 }
\end{bt}

% nguồn: MaHoaNH_TraLoiNgan_Lop10
\begin{bt}%[0H5H3-5]%[ma:0158d35d]
	Cho $4$ điểm $A$, $B$, $C$, $D$. Gọi $I$, $J$ lần lượt là trung điểm của $B C$ và $C D$. Khi đó $\overrightarrow{AB}+\overrightarrow{A I}+\overrightarrow{J A}+\overrightarrow{D A}=k \overrightarrow{D B}$. Vậy $k$ bằng  bao nhiêu?\shortans{1,5}
	\loigiai{\begin{center}
	\begin{tikzpicture}[scale=.7,>=stealth]
			\coordinate [label=left:$A$](A) at  (0,0);
			\coordinate [label=below:$B$](B) at (7,-1); 
			 \coordinate [label=above:$C$](C) at (6,5);
			\coordinate[label=above:$D$] (D) at (1,3); 
			\coordinate[label=right:$I$] (I)  at ($(C)!.5 ! (B)$);
			\coordinate[label=above:$J$] (J) at ($(C)!.5  ! (D)$);
			\foreach \diem in {A,B,C,D,I,J} \fill (\diem) circle (1pt); 
			 \draw[smooth](A)--(B)--(C)--(D)--(A)
				(I)--(A)--(J); 
		\end{tikzpicture}
\end{center}
Ta  có
\begin{eqnarray*}
	\overrightarrow{AB}+\overrightarrow{A I}+\overrightarrow{J A}+\overrightarrow{D A}&=&\overrightarrow{D A}+\overrightarrow{AB}+\overrightarrow{J A}+\overrightarrow{A I}  \\
	&= & \overrightarrow{D B}+\overrightarrow{I J}\\ 
	&= & \overrightarrow{D B}+\overrightarrow{J  I}\\
	 &=& \overrightarrow{D B}+\dfrac{1}{2} \overrightarrow{D B}\\
	 &=& \dfrac{3}{2} \overrightarrow{D B}.
\end{eqnarray*}
 }
\end{bt}

% nguồn: MaHoaNH_TraLoiNgan_Lop10
\begin{bt}%[0H5H3-5]%[ma:fa0eaf71]
	Cho hình bình hành $ABCD$ tâm $O$. Lấy các điểm $I$, $J$ sao cho $3\overrightarrow{IA}+2 \overrightarrow{IC}-2 \overrightarrow{ID}=\overrightarrow{0}$, $\overrightarrow{JA}-2 \overrightarrow{JB}+2\overrightarrow{JC}=\overrightarrow{0}$. Nếu $\overrightarrow{IJ}=k \overrightarrow{IO}$ thì  giá trị của $k$ bằng bao nhiêu?\shortans{4}
	\loigiai{\begin{center}
	\begin{tikzpicture}[scale=.7]
			\coordinate [label=above:$A$](A)  at (3,3);
			\coordinate  [label=below:$B$](B) at (0,0);
			\coordinate [label=below:$C$](C) at (8,0); 
			\coordinate[label=above:$D$] (D) at ($(C)-(B)+(A)$); 
			 \coordinate [label=below:$O$](O) at ($(A)!.5 ! (C)$);
			\coordinate [label=above left:$I$](I) at ($(A)! 2/3 !  (B)$);
			\coordinate [label=above left:$J$](J) at ($(A)!  2 ! (D)$);
			\foreach \diem in  {A,B,C,D,O,I,J} \fill (\diem) circle (1pt);
			\draw[smooth](A)--(B)--(C)--(D)--(A)--(C) 
				 (J)--(D)--(B)
				 (I)--(J);
		\end{tikzpicture}
\end{center}
 Theo đề ta có
\begin{eqnarray*}
	& & 3 \overrightarrow{IA}+2 \overrightarrow{IC}-2 \overrightarrow{ID}=\overrightarrow{0}\Leftrightarrow 3 \overrightarrow{IA}+2(\overrightarrow{IC}-\overrightarrow{ID})=\overrightarrow{0}\\ 
	&\Leftrightarrow & 3 \overrightarrow{IA}+2 \overrightarrow{DC}=\overrightarrow{0}\Leftrightarrow 3 \overrightarrow{IA}+2  \overrightarrow{AB}=\overrightarrow{0}\Leftrightarrow \overrightarrow{AI}=\dfrac{2}{3} \overrightarrow{AB}
\end{eqnarray*}
Mặt khác $\overrightarrow{JA}-2 \overrightarrow{JB}+2  \overrightarrow{JC}=\overrightarrow{0} \Leftrightarrow \overrightarrow{JA}+2\left(\overrightarrow{JC}-\overrightarrow{JB}\right)=\overrightarrow{0} \Leftrightarrow \overrightarrow{AJ}=2\overrightarrow{BC} \Leftrightarrow \overrightarrow{AJ}=2 \overrightarrow{AD}$.\\
Ta có $\overrightarrow{IO}=\overrightarrow{AO}-\overrightarrow{AI}=\dfrac{1}{2}  \overrightarrow{AC}-\dfrac{2}{3} \overrightarrow{AB}=\dfrac{1}{2}\left(\overrightarrow{AB}+\overrightarrow{AD}\right)-\dfrac{2}{3} \overrightarrow{AB}=-\dfrac{1}{6} \overrightarrow{AB}+\dfrac{1}{2} \overrightarrow{AD}$\\
và $\overrightarrow{IJ}=\overrightarrow{AJ}-\overrightarrow{AI}=2\overrightarrow{AD}-\dfrac{2}{3}  \overrightarrow{AB}=-\dfrac{2}{3}\overrightarrow{AB}+2 \overrightarrow{AD}$.\\
Do đó $\heva{&\overrightarrow{IO} = - \dfrac{1}{6} \overrightarrow{AB} + \dfrac{1}{2} \overrightarrow{AD} \\ &\overrightarrow{IJ} = - \dfrac {2}{3} \overrightarrow{AB} + 2 \overrightarrow{AD}}  \Leftrightarrow \heva{&6 \overrightarrow{IO}=-\overrightarrow{AB}+3 \overrightarrow{AD}\\&\dfrac{3}{2} \overrightarrow{IJ}=-\overrightarrow{AB}+3 \overrightarrow{A D}}\Rightarrow 6 \overrightarrow{IO}=\dfrac{3}{2} \overrightarrow{IJ} \Leftrightarrow \overrightarrow{IJ}=4 \overrightarrow{IO}$.\\Vậy $k=4$.
 }
\end{bt}


%%%%%%%%%%%%%%%% TỰ LUẬN %%%%%%%%%%%%%%%%

% nguồn: MaHoaNH_TuLuan_Lop10
\begin{bt}%[0H5H3-5]%[ma:d7b4a0f3]
	Cho $\triangle A B C$. Gọi $J$ là  điểm trên cạnh $A C$ sao cho $J A=\dfrac{2}{3} J C$. Nếu $\overrightarrow{BJ}=m\overrightarrow{BA}+n\overrightarrow{B C}$ thì giá trị $m+n$ bằng bao nhiêu?
	\loigiai{ 
		Ta có $JA=\dfrac{2}{3} J C  \Leftrightarrow 3 JA =2 JC$. Mà $\overrightarrow{JA}$ và $\overrightarrow{JC}$ ngược hướng, nên
		\begin{eqnarray*}
			3 JA =2 JC&\Leftrightarrow & 3 \overrightarrow{J  A}=-2 \overrightarrow{J C}\\
			&\Leftrightarrow &  3\left(\overrightarrow{B A}-\overrightarrow{B J}\right)+2\left(\overrightarrow{B C}-\overrightarrow{B J}\right)=\overrightarrow{0}\\
			&\Leftrightarrow  & 5 \overrightarrow{B J}=3 \overrightarrow{B A}+2 \overrightarrow{B C}\\
			&\Leftrightarrow & \overrightarrow{B  J}=\dfrac{3}{5} \overrightarrow{B A}+\dfrac{2}{5} \overrightarrow{BC}.
		\end{eqnarray*}
		 Khi đó $m+n=1$.
	 }
\end{bt}

% nguồn: MaHoaNH_TuLuan_Lop10
\begin{bt}%[0H5H3-5]%[ma:dc01da0a]
	Cho hình bình hành $ABCD$. Gọi  $E$ và $F$ là $2$ điểm thỏa $\overrightarrow{BE}=\dfrac{1}{3}$, $\overrightarrow{BC}$, $\overrightarrow{BF}=\dfrac{1}{4} \overrightarrow{BD}$. Tìm $k$ sao cho $\overrightarrow{AE}=k\overrightarrow{AF}$.
	 \loigiai{\immini{
		Ta phân tích $\overrightarrow{AE}$ và $\overrightarrow{AF}$  theo $2$ vectơ $\overrightarrow{AB}$ và $\overrightarrow{AD}$.\\
		Ta có $$\overrightarrow{AE}=\overrightarrow{AB}+\overrightarrow{BE}=\overrightarrow{AB}+\dfrac{1}{3} \overrightarrow{BC}=\overrightarrow{AB}+\dfrac{1}{3} \overrightarrow{AD}.$$ 
		Mặt khác $$\overrightarrow{AF}=\overrightarrow{AB}+\overrightarrow{BF}=\overrightarrow{AB}+\dfrac{1}{4}\left(\overrightarrow{AD}-\overrightarrow{AB}\right)=\dfrac{3}{4}  \overrightarrow{AB}+\dfrac{1}{4} \overrightarrow{AD}.$$
		 }
		{\begin{tikzpicture}[scale=0.8,font=\footnotesize,line join=round,line cap=round,>=stealth]
		\coordinate (A) at (1,3); 
		 \coordinate (B) at (6,3);
		 \coordinate (D) at (0,0);
		\coordinate (C) at  ($(B)+(D)-(A)$);
		\coordinate (E) at ($(B)!1/3!(C)$); 
		\coordinate (F)  at ($(B)!1/4!(D)$);
		 \draw (A)--(B)--(C)--(D)--cycle;
		\draw (A)--(E) 
			 (B)--(D);
		\foreach \i/\g in {A/90,B/90,C/-90,D/-90,E/0,F/90}{\draw[fill=black](\i) circle  (1pt)
			($(\i)+(\g:3mm)$)  node[scale=1]{$\i$};}
	\end{tikzpicture}
		} 
		\noindent Xét hệ $$\heva{&\overrightarrow{AE} = \overrightarrow{AB } + \dfrac {1}{3} \overrightarrow{AD} \\ &\overrightarrow{AF} = \dfrac {3}{4} \overrightarrow{AB}+ \dfrac{1}{4} \overrightarrow{AD}} \Leftrightarrow \heva{&\overrightarrow{AE}=\overrightarrow{AB}+\dfrac{1}{3} \overrightarrow{AD} \\ &\dfrac{4}{3} \overrightarrow{AF}=\overrightarrow{AB}+\dfrac{1}{3} \overrightarrow{AD}} \Leftrightarrow  \overrightarrow{AE}=\dfrac{4}{3} \overrightarrow{AF}.$$
	Vậy $k=\dfrac{4}{3}$.} 
\end{bt}
